\subsection{乘子}
\label{sec:ch3-multipliers}

给定函数 $m\in L^\infty(\mathbb R^n)$, 在 $L^2(\mathbb R^n)$ 上定义有界算子 $T_m$:
\begin{equation}
\label{eq:ch3-multiplier-definition}
 \widehat{T_mf}(\xi)=m(\xi)\widehat f(\xi).
\end{equation}
由 Plancherel 定理, 当 $f\in L^2$ 时 $T_mf$ 有良好定义, 并且
\[
 \|T_mf\|_2\le\|m\|_\infty\|f\|_2.
\]
容易看出 $T_m$ 的算子范数就是 $\|m\|_\infty$: 固定 $\varepsilon>0$, 在集合 $\{x\in\mathbb R^n:|m(x)|>\|m\|_\infty-\varepsilon\}$ 中取一个测度有限且为正的可测子集 $A$, 并令 $f$ 是满足 $\widehat f=\chi_A$ 的 $L^2$ 函数. 则
\[
 \|T_mf\|_2>(\|m\|_\infty-\varepsilon)\|f\|_2.
\]

我们称 $m$ 为算子 $T_m$ 的乘子, 尽管有时也把算子本身称为乘子. 当 $T_m$ 可以延拓为 $L^p$ 上的有界算子时, 称 $m$ 为 $L^p$ 乘子.

例如, Hilbert 变换的乘子 $m(\xi)=-i\operatorname{sgn}(\xi)$ 是 $L^p$ 乘子. 更一般地, 给定 $a,b\in\mathbb R$, $a<b$, 定义 $m_{a,b}(\xi)=\chi_{(a,b)}(\xi)$. 令 $S_{a,b}$ 为与此乘子相伴的算子:
\[
 \widehat{S_{a,b}f}(\xi)=\chi_{(a,b)}(\xi)\widehat f(\xi).
\]
它还有等价表达式
\begin{equation}
\label{eq:ch3-interval-multiplier-hilbert}
 S_{a,b}=\frac i2\bigl(M_aHM_{-a}-M_bHM_{-b}\bigr),
\end{equation}
其中 $M_a$ 是逐点乘以 $e^{2\pi iax}$ 的算子,
\[
 M_af(x)=e^{2\pi iax}f(x).
\]
为验证这一点, 注意由\eqref{eq:ch1-fourier-translation-modulation}, $iM_aHM_{-a}$ 的乘子为 $\operatorname{sgn}(\xi-a)$. 显然 $M_a$ 在 $L^p$, $1\le p\le\infty$ 上有界且范数为 $1$. 因而由\eqref{eq:ch3-interval-multiplier-hilbert} 和 $H$ 的强 $(p,p)$ 不等式, $S_{a,b}$ 在 $L^p$ 上有界.

\hypertarget{chapter-03-page-071}{}
对上述论证作显然的修改, 可知当 $a=-\infty$ 或 $b=\infty$ 时结论仍成立. 因而得到下述结果.

\begin{Proposition}
\label{prop:ch3-interval-multiplier-bound}
存在常数 $C_p$, $1<p<\infty$, 使得对所有 $-\infty\le a<b\le\infty$,
\[
 \|S_{a,b}f\|_p\le C_p\|f\|_p.
\]
\end{Proposition}

作为这一结果的一个应用, 取 $a=-R$, $b=R$. 此时 $S_{a,b}$ 就是第~\ref{chap:fourier-series-integrals}~章中引入的部分和算子 $S_R$: $S_Rf=D_R*f$, 其中 $D_R$ 是 Dirichlet 核. 因而
\[
 \|S_Rf\|_p\le C_p\|f\|_p,
\]
且常数与 $R$ 无关. 由此得到下述推论.

\begin{Corollary}
\label{cor:ch3-fourier-integral-partial-sum-convergence}
若 $f\in L^p(\mathbb R)$, $1<p<\infty$, 则
\[
 \lim_{R\to\infty}\|S_Rf-f\|_p=0.
\]
\end{Corollary}

当 $p=1$ 时没有范数收敛, 只有测度意义下的收敛:
\[
 \lim_{R\to\infty}\bigl|\{x\in\mathbb R:|S_Rf(x)-f(x)|>\varepsilon\}\bigr|=0.
\]
由这个结果和推论~\ref{cor:ch3-fourier-integral-partial-sum-convergence} 可知, 存在一个序列 $\{S_{R_k}f(x)\}$, 它对几乎所有 $x$ 收敛到 $f(x)$, 但这个序列依赖于 $f$.

给定 $L^p$ 上一族一致有界的算子, 它们的任意凸组合仍然有界. 因而由命题~\ref{prop:ch3-interval-multiplier-bound} 还可得到另一个推论.

\begin{Corollary}
\label{cor:ch3-bounded-variation-multiplier}
若 $m$ 是 $\mathbb R$ 上的有界变差函数, 则 $m$ 是 $L^p$ 乘子, $1<p<\infty$.
\end{Corollary}

\begin{Proof}
因为 $m$ 有界变差, 所以 $\lim_{t\to-\infty}m(t)$ 存在. 必要时给 $m$ 加上一个常数, 可设该极限为零. 还可将 $m$ 规范化为在每个 $x\in\mathbb R$ 处右连续. 令 $dm$ 为与 $m$ 相伴的 Lebesgue--Stieltjes 测度. 则
\[
 m(\xi)=\int_{-\infty}^{\xi}dm(t)
 =\int_{\mathbb R}\chi_{(-\infty,\xi)}(t)\,dm(t)
 =\int_{\mathbb R}\chi_{(t,\infty)}(\xi)\,dm(t).
\]
因此
\[
 \widehat{T_mf}(\xi)
 =\int_{\mathbb R}\chi_{(t,\infty)}(\xi)\widehat f(\xi)\,dm(t),
\]

\hypertarget{chapter-03-page-072}{}
从而
\[
 T_mf(x)=\int_{\mathbb R}S_{t,\infty}f(x)\,dm(t).
\]
由 Minkowski 不等式,
\[
 \|T_mf\|_p
 \le\int_{\mathbb R}\|S_{t,\infty}f\|_p\,|dm|(t)
 \le C_p\|f\|_p\int_{\mathbb R}|dm|(t).
\]
$|dm|$ 的积分就是 $m$ 的全变差, 按假设它是有限的.
\end{Proof}

给定一个乘子, 可以通过平移、伸缩和旋转构造其他乘子.

\begin{Proposition}
\label{prop:ch3-multiplier-transformations}
若 $m$ 是 $L^p(\mathbb R^n)$ 乘子, 则由
\[
 m(\xi+a),\quad a\in\mathbb R^n;
 \qquad m(\lambda\xi),\quad\lambda>0;
 \qquad m(\rho\xi),\quad\rho\in O(n)
\]
定义的函数都是 $L^p$ 上有界算子的乘子, 且算子范数均与 $T_m$ 相同.
\end{Proposition}

这一结论由 Fourier 变换的性质\eqref{eq:ch1-fourier-translation-modulation}、\eqref{eq:ch1-fourier-orthogonal-transform} 与\eqref{eq:ch1-fourier-dilation} 立即得到.

若 $m$ 是 $L^p(\mathbb R)$ 乘子, 则 $\mathbb R^n$ 上的函数 $\widetilde m(\xi)=m(\xi_1)$ 是 $L^p(\mathbb R^n)$ 乘子. 事实上, 若 $T_m$ 是与 $m$ 相伴的一维算子, 则对定义在 $\mathbb R^n$ 上的 $f$,
\[
 T_{\widetilde m}f(x)=T_mf(\,cdot\,,x_2,\ldots,x_n)(x_1).
\]
由 Fubini 定理和 $T_m$ 在 $L^p(\mathbb R)$ 上的有界性,
\[
\begin{aligned}
 \int_{\mathbb R^n}|T_{\widetilde m}f(x)|^p\,dx
 &=\int_{\mathbb R^{n-1}}
   \left(\int_{\mathbb R}|T_mf(\,cdot\,,x_2,\ldots,x_n)(x_1)|^p\,dx_1\right)
   dx_2\cdots dx_n\\
 &\le C\int_{\mathbb R^{n-1}}
   \left(\int_{\mathbb R}|f(x_1,\ldots,x_n)|^p\,dx_1\right)
   dx_2\cdots dx_n.
\end{aligned}
\]

取 $m=\chi_{(0,\infty)}$. 由命题~\ref{prop:ch3-interval-multiplier-bound}, $m$ 是 $L^p(\mathbb R)$ 乘子, $1<p<\infty$. 因而由前面的论证, 半空间 $\{\xi\in\mathbb R^n:\xi_1>0\}$ 的特征函数是 $L^p(\mathbb R^n)$ 乘子. 再由命题~\ref{prop:ch3-multiplier-transformations}, 任意半空间的特征函数也具有同样性质, 因为任意半空间都可由上述半空间经旋转和平移得到. 一个有 $N$ 个面的凸多面体的特征函数可写成 $N$ 个半空间特征函数的乘积, 所以它也是 $L^p(\mathbb R^n)$ 乘子, $1<p<\infty$. 由此得到下述推论.

\begin{Corollary}
\label{cor:ch3-convex-polyhedron-summability}
若 $P\subset\mathbb R^n$ 是包含原点的凸多面体, 则
\[
 \lim_{\lambda\to\infty}\|S_{\lambda P}f-f\|_p=0,
 \qquad 1<p<\infty,
\]
其中 $S_{\lambda P}$ 是以 $\lambda P=\{\lambda x:x\in P\}$ 的特征函数为乘子的算子.
\end{Corollary}

\hypertarget{chapter-03-page-073}{}
根据上述观察, 下述事实可能颇令人惊讶: 当 $n>1$ 时, 以原点为中心的球的特征函数在 $p\ne2$ 时不是 $L^p(\mathbb R^n)$ 上的有界乘子. 我们将在\MBUnresolvedReference{chapter}{第 8 章}中考察这一事实以及 $\mathbb R^n$ 上乘子的相关结果(另见下文第~\ref{subsec:ch3-fourier-series-multipliers}~节).
